% Compile from this directory with: xelatex why_three_qiskit6.tex
\documentclass[11pt,a4paper]{article}
\usepackage{kotex}
\usepackage{amsmath,amssymb,booktabs,geometry,hyperref}
\geometry{margin=25mm}
\hypersetup{colorlinks=true,urlcolor=blue}
\title{I01 Qiskit 실험 6: 균일 기준점의 시간 에일리어싱과 독립 검증점}
\author{RELATIONAL ATLAS --- 유한 시계의 위상 기준 샘플링}
\date{2026년 10월 3일}

\begin{document}
\maketitle

\begin{abstract}
기준 측정을 더 자주 한다는 것만으로 과학 측정 시점의 위상 영점을
복원할 수 있는지 시험했다. 다섯 개의 균일 기준점 모두에서 값이 0인
숨은 진동을 주입하고 그 사이에서 과학 측정을 한다. 이 구성에서
다섯 점은 위상 진동의 존재 자체를 놓치며, 비제약 과학계의 공통 위상과
미지 진동 진폭은 정확히 혼동되어 Fisher rank 3이다. 격자 사이의
독립 검증점은 그 진동을 드러내고, 과학 측정 시점의 기준점을 포함한
아홉 점 분석은 위상 추정을 복원한다. Aer의 유한-shots 결과에서도
다섯 점 보간은 참값 $0.300$ rad를 $0.115$ rad로 오독했고, 직접
기준점은 $0.294$ rad로 복원했다. 유한 수의 기준점이 모든 가능한
드리프트를 보장할 수는 없으며, 어떤 시간적 대역 또는 매끄러움의
가정을 쓰는지를 명시해야 한다.
\end{abstract}

\section{실험 설계와 숨은 모드}
이전과 같은 비제약/$P=0$ 4축 토이 상태, 5준위 에너지 매칭 시계,
9개 위상 표지, 동일한 복귀 검출기를 사용한다. 한 기준 측정 구간의
정규화 시간 $u\in[0,1]$에서 드리프트를
\begin{equation}
 \delta(u)=\beta_0+\beta_1u+A\sin(4\pi u),\qquad
 \beta_0=-0.150,\quad\beta_1=0.500,\quad A=0.180\ \mathrm{rad}
 \label{eq:drift}
\end{equation}
로 주입한다. 과학 측정은 $u_s=3/8$에 두고 실제 과학 공통 위상은
$\alpha=0.300$ rad로 한다. 다섯 균일 기준점
$\{0,1/4,1/2,3/4,1\}$에서 $\sin(4\pi u)=0$이므로 주입된 진동을
\emph{아무 정밀도로 측정해도} 이 다섯 점만으로는 볼 수 없다.
반면 $u_s=3/8$에서는 사인값이 $-1$이다. 실제 과학 시점의 영점은
$\delta(u_s)=-0.1425$, 과학 측정에 들어가는 합성 위상은
$\alpha+\delta(u_s)=0.1575$ rad다. 가까운 기준점 $1/4,1/2$의
직선 보간은 영점을 $0.0375$ rad로 예측해 $\alpha$를 $0.120$ rad로
오독한다.

아홉 점 격자 $u=j/8$을 별도로 기록한다. 이 중 짝수 $j$가 원래
다섯 점이며 홀수 $j=1,3,5,7$는 사전에 정한 독립 검증점이다.
$j=3$은 과학 시점 자체의 알려진 위상 기준점이므로, 그 직전/직후에
위상 영점이 안정적이라는 자원 가정 아래 모델 형태를 몰라도
$\delta(u_s)$를 직접 추정할 수 있다. $j=1$ 등 다른 검증점은
기존 다섯 점의 보간이 틀린다는 경고 신호다.

\section{정확 Fisher 정보와 조건부 rank}
숨은 사인 진폭 $A$를 \emph{가능한 nuisance 모드}로만 두고
$(\boldsymbol\theta,\beta_0,\beta_1,A)$의 결합 Fisher 행렬에서
세 nuisance를 제거했다. 다섯 점에서는 $A$의 기준 측정 미분이
모두 0이고 과학 시점에서는 $-1$이다. 따라서
\begin{equation}
 (\boldsymbol\theta,A)\mapsto
 (\boldsymbol\theta+c\boldsymbol1,A+c)
\end{equation}
가 과학계와 다섯 기준점의 확률을 모두 보존한다. 그 결과는 아래와
같다. 다섯 시계 위상 $\tau/\pi=0,1/4,1/2,3/4,1$에서 동일하다.
\begin{center}
\begin{tabular}{@{}lccc@{}}
\toprule
과학계 & 균일 5점 & 5점+$u=1/8$ 검증점 & 균일 9점 \\
\midrule
비제약 & 3 & 4 & 4 \\
$P=0$ & 3 & 3 & 3 \\
\bottomrule
\end{tabular}
\end{center}
중간 열에서 rank 4의 회복은 \emph{사인 모드라는 후보 모형을 알고}
$u=1/8$의 측정으로 그 진폭을 추정하는 조건부 결과다. 사인 형태를
전혀 가정하지 않은 채 단 하나의 다른 시점 기준값으로 임의의
$u_s$ 영점을 결정할 수는 없다. 마지막 열의 유한-shots 분석은 사인
함수를 적합하지 않고 $u_s$의 기준 측정값을 직접 사용한다.

\section{Aer 유한-shots: 보간 편향과 홀드아웃 잔차}
시계 위상 $\tau=0,\pi/2$ 각각에 비제약 과학 36회로와 기준 9점
각 36회로를 배치해 720회로를 실행했다. $P=0$의 실제/무진동
공통 위상 대조 144회로도 추가했다. 전체 864회로에 각각 4096
shots, 최대 CPU 병렬 6을 썼다. 다섯 점 분석은 같은 비제약 raw
counts 중 과학과 짝수 기준 5점만 사용한다. 시계 후선택률은
비제약 $0.20000$, $P=0$ $0.19918$이었다.
\begin{center}
\begin{tabular}{@{}ccccc@{}}
\toprule
$\tau/\pi$ & 5점 보간 $\hat\alpha$ & 9점 직접 기준 $\hat\alpha$
& $u=1/8$ 잔차 & $u=3/8$ 잔차 \\
\midrule
0 & $0.1150\ [0.1006,0.1294]$ & $0.2940\ [0.2781,0.3099]$
& $+0.1860\ [0.1722,0.1998]$ & $-0.1790\ [-0.1921,-0.1659]$ \\
1/2 & $0.1348\ [0.0834,0.1861]$ & $0.3230\ [0.2636,0.3824]$
& $+0.1923\ [0.1407,0.2438]$ & $-0.1883\ [-0.2398,-0.1367]$ \\
\bottomrule
\end{tabular}
\end{center}
단위는 rad, 대괄호는 세그먼트별 국소 Fisher 표준오차를 합친 근사
95\% 구간이다. 잔차는 홀수 기준점의 적합 위상에서 양옆 짝수
기준점의 평균을 뺀 값이다. $u=1/8$ 검증점의 잔차는 과학 데이터를
쓰지 않고도 5점 선형 보간의 실패를 발견한다. $u=3/8$ 기준점은
그 실패를 과학 위상 보정까지 연결한다. $\tau=\pi/2$의 오차구간은
조건부 신호 감소 때문에 더 넓다.

$P=0$의 정확한 공통 위상 불변성은 운동량 합 제약으로 검증했다.
첫 seed의 두 시계 위상 재배분 검정 $p$는 $0.765,0.040$이었다.
뒤의 $0.040$을 물리적 변화로 해석하지 않는다. 독립 seed 15개를
추가해 전체 16 seed$\times$2 시계 위상, 32개 검정 가운데
5개가 $p\le0.05$였다. 명목 기대값 1.6개보다 많지만 이 작은
ensemble만으로 Aer 샘플러나 검정의 체계적 편차를 확정할 수도 없다.
검정식 자체를 별도로 점검하려고 정확 $P=0$ 이항확률에서 각 위상
1000쌍의 독립 표본을 만들고 각 표본에 199회 조건부 재배분을 적용했다.
두 시계 위상 모두 53/1000이 $p\le0.05$로, 명목 5\%에 근접했다.
따라서 Aer 대조군의 작은 $p$값은 현 단계에서 \emph{미해결 유한표본
품질 점검 사항}으로 남기며, 정확 불변성이나 주된 보간 개입의
증거와 혼동하지 않는다. 모든 seed와 raw counts를 JSON에 보존했다.

\section{해석과 범위}
기준점의 개수만으로 시간 해상도를 정의할 수 없다. 기준점 배치가
변동 모드와 맞물리면 에일리어싱이 생긴다. 과학 시점의 기준 측정은
이 사례를 해결하지만, 실제 장치에서 과학/기준 샷 사이에 더 빠른
드리프트가 있다면 다시 실패할 수 있다. 유한한 기준점 집합이
과학 시점을 포함하지 않는다면, 그 기준점에서 0이고 과학 시점에서
양수인 매끄러운 함수를 만들 수 있으므로 \emph{임의의} 드리프트에
대한 보장은 원리적으로 불가능하다. 따라서 앞으로는 시간 대역폭,
매끄러움, 기준 측정 간격, 그리고 샷 예산을 함께 명시해야 한다.

숨은 사인 모드는 이 문제를 드러내기 위해 설계한 반례이며 실제
양자 장치의 잡음 법칙으로 주장하지 않는다. 완벽한 시계, $P=0$
상태, 알려진 위상의 기준 측정도 모두 준비 자원이다. 이 결과는
자발적인 세 공간 차원 선택이 아니다. QPU는 사용하지 않았다.

\paragraph{재현성.} 코드
\texttt{../src/relational\_atlas\_i01/phase\_reference/qiskit\_clock\_alias\_audit.py}, 회귀 검사
\texttt{../tests/test\_qiskit\_clock\_alias\_audit.py}, 원자료
\texttt{../results/qiskit\_clock\_alias\_audit.json}.
시작 seed 20261008, Qiskit 2.4.2, Qiskit Aer 0.17.2.
Windows Aer의 합성 \texttt{multiplexer} 접근 위반을 피하려고
준비 회로를 $u3,cx$로 분해했다. TeX 엔진이 없어 PDF 렌더링은
아직 검증하지 못했다.
\end{document}
